\documentclass{article}
\usepackage{amsmath,amssymb,hyperref}
\begin{document}
\section*{The union problem for hyperfinite Borel equivalence relations}
\subsection{Definitions and mathematical statement}
A standard Borel space is a set equipped with the Borel structure of some
Polish topology. On such a space $X$, a Borel equivalence relation is a Borel
subset of $X\times X$ that is an equivalence relation. It is finite when all
its classes are finite, with no common bound on their sizes required.
It is hyperfinite when it has an exhaustion
\[
 R=\bigcup_{k\in\mathbb N}R_k,\qquad R_k\subseteq R_{k+1},
\]
by finite Borel equivalence relations on $X$.

Suppose $F_1\subseteq F_2\subseteq\cdots$ are hyperfinite Borel equivalence
relations on the same standard Borel space $X$. Must
\[
 E=\bigcup_{n\in\mathbb N}F_n
\]
be hyperfinite? The conclusion asks for one increasing exhaustion of $E$
by finite Borel equivalence relations. All assertions concern the whole
space, without discarding a null set or a meagre set.
\subsection{Short English statement}
Does taking an increasing countable union preserve hyperfiniteness for Borel equivalence relations?
\subsection{Sources}
S. Gao and M. Xiao, \emph{An order analysis of hyperfinite Borel equivalence
relations}, arXiv:2404.17516v2, Problem 2.9 and Section 2,
\url{https://arxiv.org/html/2404.17516v2}.
Adapted under CC BY 4.0, \url{https://creativecommons.org/licenses/by/4.0/}.
The general question is also stated as open in the introduction to
\emph{Hyperfiniteness of bounded-to-one actions of commutative monoids},
arXiv:2608.18439v2, \url{https://arxiv.org/html/2608.18439v2}.
\end{document}
